\section{Related Work}
\label{sec:related}

\paragraph{RAG systems and retrieval-based QA.}
The foundational RAG framework of \citet{lewis2020rag} showed that coupling a parametric generator with a nonparametric retriever substantially improves performance on knowledge-intensive tasks. Retrieval-augmented pretraining, explored in REALM~\citep{guu2020realm} and scaled dramatically in RETRO~\citep{borgeaud2022retro}, established that grounding language model training in retrieved documents improves factual accuracy across very large corpora. Dense retrieval has since split into two dominant paradigms---dual-encoder models like DPR~\citep{karpukhin2020dpr} and late-interaction models like ColBERT~\citep{khattab2020colbert}---while learned sparse representations from SPLADE~\citep{formal2021splade} offer a middle ground that preserves BM25-style interpretability. A systematic comparison of these families by \citet{luan2021sparse} shows that no single paradigm dominates across all retrieval settings. More recent work has focused on the generation side: Self-RAG~\citep{asai2024selfrag} introduces reflection tokens that let the generator selectively retrieve and critique evidence, and \citet{pradeep2023rankgpt} demonstrate that instruction-tuned LLMs can serve as effective passage re-rankers. The BERGEN benchmarking library~\citep{thakur2024bergen} offers a systematic framework for comparing these pipelines, but evaluates against fixed corpora. None of this work measures how retrieval quality degrades as indexed documents age relative to a rolling deployment window---the gap \method{} fills.

\paragraph{Benchmark validity, contamination, and temporal robustness.}
The fundamental threat to any static benchmark is that large-scale pretraining corpora likely overlap with evaluation sets~\citep{magar2022data,golchin2023time}. GPT-4's web-scale training~\citep{openai2024gpt4} has raised the stakes: the probability that a benchmark question appears verbatim in pretraining data has grown from a theoretical concern to a practical one. The HELM framework~\citep{liang2022helm} exposed how aggregate leaderboard scores can obscure this variation across dozens of evaluation scenarios. Dynabench~\citep{kiela2021dynabench} and LiveBench~\citep{white2024livebench} attack the problem differently---through adversarial human collection and evaluation against recently published problems, respectively---though neither addresses temporal drift in retrieval corpora specifically. The data-centric survey of \citet{longpre2023pretraining} identifies benchmark contamination as a systemic threat to reproducible evaluation across NLP.

The temporal dimension of this problem has received growing attention. FreshQA~\citep{vu2023freshqa} showed that closed-book model accuracy on world-knowledge questions decays measurably within months of a knowledge cutoff. The ``temporal half-life'' concept from \citet{kasner2024beyond} applies directly here: QA accuracy on news-centric topics varies by domain as a function of how quickly the underlying world changes. Yet both FreshQA and the half-life analysis focus on closed-book settings. \method{} extends this line to retrieval-grounded systems, where the freshness of the \emph{corpus} is as important as the freshness of the \emph{model}---and where citation precision offers an independent signal of staleness that answer accuracy alone misses.

\paragraph{Citation faithfulness and factual grounding in generation.}
Whether generated text is faithfully grounded in retrieved evidence has emerged as a distinct evaluation challenge. \citet{gao2023enabling} proposed using NLI to verify whether each generated sentence is supported by at least one cited passage, achieving reasonable agreement with human judges. Concurrent work on GopherCite~\citep{menick2022gophercite} and WebGPT~\citep{nakano2021webgpt} demonstrated that explicit grounding constraints substantially reduce hallucination rates. FActScoring~\citep{min2023factscore} pushed this further by decomposing long-form generation into atomic claims and assessing each against a reference corpus, revealing that surface fluency and factual support diverge significantly. RAGAS~\citep{es2023ragas} packages faithfulness, answer relevance, and context precision into a unified suite, while MiniCheck~\citep{tang2024minicheck} shows that lightweight models can approximate NLI-based citation checking at substantially reduced cost.

What this line of work lacks, by construction, is a temporal dimension: existing citation evaluation frameworks treat the retrieval corpus as fixed. A key finding of \method{} is that citation faithfulness degrades 2.3$\times$ faster than answer fluency on fresh slices---an effect invisible to any evaluation framework that does not distinguish parametric recall from retrieval-grounded generation.

\paragraph{Temporal evaluation in NLP.}
Beyond RAG, temporal robustness has been studied in information extraction~\citep{loureiro2022timelms}, relation classification~\citep{dhingra2022time}, and event detection~\citep{cohen2021temporal}. \citet{chen2022temporalqa} showed that accuracy for both closed-book models and retrieval-augmented systems degrades as the temporal gap between question creation and model knowledge grows. Domain shift under temporal distribution change is particularly stark in biomedical NLP, where \citet{muller2022covid} demonstrated that models fine-tuned on pre-pandemic data underperform substantially on COVID-era documents. \citet{shi2024ircot} add another wrinkle: language models are easily distracted by irrelevant context, a failure mode that compounds on fresh corpora where topically adjacent but temporally stale passages frequently appear in the top-$k$ results.

The BEIR benchmark~\citep{thakur2021beir} remains the standard heterogeneous retrieval evaluation suite, yet its splits are frozen at construction time---precisely the limitation \method{} addresses. KILT~\citep{petroni2021kilt} similarly aligns multiple knowledge-intensive tasks to a shared Wikipedia snapshot; \method{} adopts the same alignment principle but refreshes the snapshot monthly. TempLAMA~\citep{dhingra2022time} is the closest concurrent work, probing parametric temporal knowledge in language models. Unlike TempLAMA, \method{} targets system-level behavior rather than model internals, and includes retrieval as a first-class component of the evaluation.
